\chapter{Troubleshooting and FAQ}\label{ch:troubleshooting}

\section{Guard failures (plan stops with ``Resource precondition failed'')}
The message names the guard. Fix the model; never disable the guard to make a plan pass.
{\small
\begin{longtable}{@{}>{\raggedright\arraybackslash}p{0.36\linewidth}>{\raggedright\arraybackslash}p{0.6\linewidth}@{}}
\toprule
\textbf{Message starts with} & \textbf{Cause and fix} \\ \midrule\endhead
Duplicate role name & The same role key appears in two \texttt{roles/*.yaml} files. Keep one. \\
Object names must be absolute & A name lacks the leading \texttt{/} or ends with \texttt{/}. \\
Role names must start with / Auth method names must start with & Naming policy (\texttt{var.naming}). Rename, or set the prefix to \texttt{""} to disable the check. \\
Unknown auth method type & Supported in the model: api\_key, oidc, saml, oauth2, aws\_iam, azure\_ad, k8s, universal\_identity. Others: manage with \texttt{export --hcl}. \\
Unsupported rule type & Use item-, target-, role-, auth-method- or sra-rule; \texttt{ara-rule} needs \texttt{allow\_ara\_rules = true} after you verified it. \\
Invalid capability for the rule type & e.g.\ \texttt{allow\_access} on an item-rule. See the list in the message. \\
A deny rule must contain only 'deny' & Split it: one rule with \texttt{[deny]}, another with the allowed capabilities. \\
approval\_authority cannot be combined & Put approvers in a separate rule or role. \\
Root-scope rule & \texttt{/*} with more than list/deny. Narrow the path, or set \texttt{allow\_root\_scope: true} with a documented reason. \\
Invalid management-plane value & e.g.\ \texttt{audit\_access: some}. Allowed values are in the model reference. \\
Association references an auth method that is not in & Add the auth method to the model, or set \texttt{allow\_external\_auth\_methods: true} on the role. \\
Auth method has force\_sub\_claims = true & Every association to it needs \texttt{sub\_claims}. \\
PKI issuer: \ldots & Exactly one of \texttt{signer\_key\_name} / \texttt{ca\_target}; \texttt{auto\_renew} needs \texttt{destination\_path}; CRL/OCSP/ACME/public CA need \texttt{gw\_cluster\_url}. \\
custom\_fields.object\_type & STATIC\_SECRET, DYNAMIC\_SECRET or ROTATED\_SECRET. \\
event\_forwarders: type must be & Type or event name not valid; \texttt{watchover forwarder} validates both. \\
event\_forwarders: runner\_type periodic & Add \texttt{every} (hours). \\
event\_forwarders: slack and teams need url\_var & The webhook URL is a secret: add \texttt{url\_var} and the key to \texttt{TF\_VAR\_event\_forwarder\_secrets}. \\
event\_forwarders of type teams require & Add \texttt{gateway\_event\_source\_locations} and \texttt{runner\_type}. \\
model/gateway.yaml defines gateway objects & Set \texttt{gateway\_api} (CI variable \texttt{AKEYLESS\_GATEWAY\_API}). \\
gateway allowed\_access: \ldots & Unknown permission, or both/neither of \texttt{auth\_method} and \texttt{access\_id}. \\
gateway log\_forwarding.type & syslog or splunk; other targets need an extra resource block. \\
\bottomrule
\end{longtable}}

\section{Pipeline and Terraform}
\begin{description}[style=nextline,leftmargin=1.5em]
\item[\texttt{argument "jwt" is required}] \texttt{TF\_VAR\_pipeline\_jwt} is not exported in the job.
  The provider's environment fallback is not enough; the variable must be set.
\item[Authentication fails (401) in CI] The token's claims do not match the role binding: wrong
  repository/project path, a branch other than the protected default branch for apply, or an
  audience other than \texttt{akeyless}. Compare the token's claims with \texttt{bootstrap.tfvars}.
\item[Plan on a merge request fails with an authorization error] The plan role lacks a read right
  (\textbf{VERIFY} item). Extend \texttt{bootstrap/} for the object family named in the error.
\item[Plan wants to \emph{replace} a role or an auth method] Its name (YAML key) changed. Names
  are ForceNew; see ``Rename a role or an auth method in Chapter~\ref{ch:workflows}''.
\item[Plan shows changes right after an apply] A value is normalised by Akeyless (case, ordering,
  defaults). Write it in the model the way Akeyless returns it; \texttt{export} shows that form.
\item[Nightly drift job fails] Someone changed Akeyless outside Git; see ``Handle a drift finding''.
\item[OpenTofu panics in a test with import blocks] Known OpenTofu behaviour with mocked providers;
  the kit validates import blocks but plans the imported model without them.
\end{description}

\section{Kit commands}
\begin{description}[style=nextline,leftmargin=1.5em]
\item[\texttt{watchover test}: \texttt{test\_usage\_guide\_is\_generated\_and\_current} fails]
  A command or option changed. Run \texttt{scripts/watchover help --markdown > docs/usage-guide.md}
  and \texttt{python3 docs/manual/build.py}.
\item[\texttt{verify}: ``describe-permissions \ldots failed''] Check that the profile works
  (\texttt{akeyless list-roles --profile <p>}). If the call fails only for paths the persona cannot
  access, your CLI reports ``no access'' as an error: use \texttt{--error-means-no-access}.
\item[\texttt{verify}: everything skipped] No persona has a \texttt{profile:} key.
\item[\texttt{export}: a type is missing] \texttt{gateway} is not in the default \texttt{--types};
  event forwarders cannot be listed by Akeyless (use \texttt{watchover forwarder}).
\item[\texttt{export --transport api}: 401] The token expired or is missing: pass \texttt{--token}
  or set \texttt{AKEYLESS\_TOKEN}; tokens from \texttt{watchover token} are short-lived.
\item[\texttt{token}: ``authentication failed''] Wrong access ID, an expired JWT, or claims that do
  not match the CI auth method's bound claims.
\item[\texttt{account-settings} locked you out] A \texttt{bound-ips} value excluded the CI runners
  or your network. Fix it as an administrator from an allowed network; always \texttt{--dry-run} first.
\item[\texttt{forwarder --write}: ``already exists''] The name is in \texttt{event\_forwarders};
  edit it instead, or choose another name.
\end{description}

\section{Frequently asked questions}
\begin{description}[style=nextline,leftmargin=1.5em]
\item[Can I edit roles in the console?] You can, but the next apply reverts it. Change Git.
\item[Why YAML and not plain HCL?] Reviewers read access, not Terraform; the guards and the
  simulator need a uniform model; the HCL stays generic and tested.
\item[Can I use plain HCL for objects outside the model?] Yes: \texttt{export --hcl --generate}
  writes plain resources; keep them next to the model in the same root module.
\item[API or CLI?] Both reach the same operations (Chapter~\ref{doc-api-vs-cli}). Use the CLI on
  workstations and the REST API in locked-down runners.
\item[Several Gateway clusters?] One state or workspace per cluster, each with its own
  \texttt{gateway\_api} and \texttt{model/gateway.yaml}.
\item[Does the skill need access to my account?] No. It produces files and commands; you run them.
\end{description}
